% EDIT (2026-10-29): rewritten to the author's outline ("showing how prompt design affects the
% outcome ... we show 1) the impact of guardrails 2) the importance of profiling data 3) ..."). Every
% number is one already reported in the body; the per-variant first-attempt counts are read off
% Table~\ref{tab:ablation}.
We present an empirical study of how prompt design affects what a large language model does when it
is asked to make GPU code faster, and of how far the gains it finds survive on production scientific
applications. Our framework, OPTIMAS, closes the loop from profiling to prompting, generation, build,
correctness check and timed measurement, with up to three repair iterations driven by compiler and
runtime feedback. We vary the prompt one component at a time and measure the outcome on HeCBench
kernels and on two production applications, AutoDock-GPU and Cholla, each checked against its own
scientific output. We show (1) the impact of guardrails: which instructions the prompt carries decides
how often a model returns a valid candidate at the first attempt (between $8$ and $11$ of $12$ runs,
depending on which component is removed) and how slow its worst result is; and without an explicit
semantic-preservation rule both models ``optimize'' a benchmark by deleting the computation it exists
to measure, reporting $1504\times$ and $2809\times$ while the benchmark's own check still passes;
(2) the importance of profiling data: removing it from an otherwise identical prompt costs
$7.7$--$11.5\%$ on a 21k-line application, while on single-kernel benchmarks it moves individual
results in both directions with no net gain; (3) the limits of interpreting that data for the model:
converting profiler metrics into a checklist of named remedies produces both our largest gain
($1.00\times \to 1.38\times$) and our largest loss ($0.84\times \to 0.35\times$), so its effect is a
wider spread of outcomes rather than a better average; and (4) what the models can and cannot do:
they reliably find local, exact transformations---on Cholla all four outputs under the baseline
prompt ($1.48$--$1.58\times$) beat our hand-tuned reference ($1.47\times$)---and reliably miss
cross-cutting restructuring, leaving unfound on AutoDock-GPU the reduction fusion that is worth
$1.53\times$ by hand. We also report the problems that appear only on real code and the measurement
practice needed to report any of this credibly on shared GPUs.

% ---------------------------------------------------------------------------
% ORIGINAL ABSTRACT (kept for reference). It describes three prompt designs
% (loose / constrained / agentic) and a one-shot agentic result; the version above
% was rewritten to match the experiments currently in the paper (six-variant
% ablation, semantic-preservation constraint, two applications, profile-data
% ablation). Restore or merge if the agentic-prompt experiment is added.
%
% We present an empirical study of how progressively adding task-specific guardrails to prompts
% impacts large language models' ability to generate code optimization suggestions for HPC workloads.
% Across three prompt designs, we observe that tighter guardrails systematically shift model behavior
% from aggressive but low-utility transformations toward conservative, high-utility optimizations that
% compile, preserve numerical correctness, and yield measurable performance gains. The first prompt,
% which loosely describes the optimization task, produces many syntactically plausible but
% non-compiling or numerically incorrect suggestions and often requires up to ten
% optimization-feedback iterations to reach a viable configuration. The second prompt introduces
% explicit constraints on correctness, safety, and resource usage, reducing over-aggressive
% transformations and increasing the fraction of optimizations that lead to positive speedups without
% degrading numerical accuracy. Finally, our most constrained "agentic" prompt instantiates a
% closed-loop optimization routine in which the LLM autonomously proposes an optimization, checks
% compilation and numerical correctness against a reference, and revises its plan before returning a
% single final suggestion, achieving one-shot successful optimization on our benchmarks. We evaluate
% these prompts on diverse HPC kernels and systems, and we quantify their effects along three axes:
% performance gain, correctness, and optimization aggressiveness (e.g., code churn, magnitude of
% configuration change). Our results show that guardrail design is not merely a safety or policy
% concern, but a central lever for turning general-purpose code LLMs into reliable performance
% co-designers for scientific computing.
% ---------------------------------------------------------------------------
